\documentclass[10pt]{extarticle}
\usepackage[a4paper,margin=1.5cm]{geometry}
\usepackage[T1]{fontenc}
\usepackage{amsmath,amssymb}
\usepackage{graphicx}
\usepackage[export]{adjustbox}
\usepackage{multicol}
\usepackage{enumitem}
\usepackage{booktabs,array,tabularx}
\setlength{\columnsep}{0.7cm}
\setlength{\parindent}{0pt}
\setlength{\parskip}{3pt}
\pagestyle{empty}
\makeatletter
% figure registry (figdefs.tex is generated by the extract step)
\newcommand{\qfigdef}[4]{\@namedef{qf@file@#1}{#2}\@namedef{qf@w@#1}{#3}\@namedef{qf@wide@#1}{#4}}
\newcommand{\qf@check}[1]{\@ifundefined{qf@file@#1}{\PackageError{score2}{Unknown figure id #1}{}}{}}
\newcommand{\qf@img}[4][]{\includegraphics[width=#2\dimexpr\@nameuse{qf@w@#4}pt\relax,max width=#3,max height=0.45\textheight,#1]{\@nameuse{qf@file@#4}}}
% Full-width questions: finalize wraps every question that contains a wide figure in \qfullbegin...\qfullend,
% so the WHOLE question (stem, figures, options) is set across the page in reading order, then the two columns
% resume. The question's head - its number, stem and first figure (\qfullsplit marks the end of the head) -
% is kept on one page, so a number is never left at a page foot with its figure overleaf; the rest may break
% across pages normally, so no big gaps are left. Without a \qfullsplit, the whole question is kept together
% if it is under 45% of a page. Never write these yourself.
\newif\ifqf@full
\newif\ifqf@boxed
\newsavebox{\qf@box}
\newcommand{\qf@open}{\global\qf@fulltrue\par\global\qf@boxedtrue
  \setbox\qf@box\vbox\bgroup\hsize\textwidth\linewidth\textwidth}
\newcommand{\qf@emit}[1]{\par\egroup\global\qf@boxedfalse
  \ifdim\dimexpr\ht\qf@box+\dp\qf@box\relax<#1\textheight\noindent\usebox{\qf@box}\par\else\unvbox\qf@box\fi}
\newcommand{\qfullsplit}{\ifqf@boxed\qf@emit{0.85}\fi}
\newcommand{\qf@close}{\ifqf@boxed\qf@emit{0.45}\fi\global\qf@fullfalse}
\newcommand{\qfullbegin}{\end{multicols}\qf@open}
\newcommand{\qfullend}{\qf@close\begin{multicols}{2}}
\newcommand{\qfullstart}{\qf@open}
\newcommand{\qfulldone}{\qf@close}
% \qfig[scale]{id}: block figure, centred, directly where it is written (size set by the tool).
% Fallback if a wide figure is not inside a full-width question: break out of the columns in place - no floats.
\newcommand{\qfig}[2][1]{\qf@check{#2}%
  \ifqf@full\qf@blk{#1}{#2}\else
  \ifnum\@nameuse{qf@wide@#2}=1 \expandafter\@firstoftwo\else\expandafter\@secondoftwo\fi
  {\end{multicols}\begin{center}\qf@img{#1}{\textwidth}{#2}\end{center}\begin{multicols}{2}}%
  {\qf@blk{#1}{#2}}\fi}
\newcommand{\qf@blk}[2]{\par\smallskip{\centering\qf@img{#1}{\linewidth}{#2}\par}\smallskip}
% \qfigi[scale]{id}: inline figure (inside option lists / tables / minipages); never breaks columns.
\newcommand{\qfigi}[2][1]{\qf@check{#2}\qf@img[valign=c]{#1}{\linewidth}{#2}}
\makeatother
% structure macros - lint reads these to check completeness and figure placement
\newcommand{\papersec}[1]{\par\bigskip{\centering\Large\bfseries PAPER-#1\par}\medskip}
\newcommand{\partsec}[2]{\par\medskip\noindent{\large\bfseries PART-#1: #2}\par\smallskip}
\newcommand{\sect}[1]{\par\smallskip\noindent\textbf{#1}\par}
\newcommand{\qn}[1]{\par\smallskip\noindent\textbf{#1.}~}
\newcommand{\paraq}[2]{\par\smallskip\noindent\textbf{Paragraph for Questions #1 and #2}\par}
\newenvironment{qoptions}{\begin{itemize}[leftmargin=1.8em,itemsep=1pt,topsep=2pt,label={}]}{\end{itemize}}
\newcommand{\opt}[2]{\item[(#1)] #2}
% figure registry (generated by score2_tools.py finalize)
\qfigdef{p1_phy_s1_q01_1}{Score Advanced - Test 11 (09 May 2023)-img/p1_phy_s1_q01_1.png}{184.4}{0}
\qfigdef{p1_phy_s1_q02_1}{Score Advanced - Test 11 (09 May 2023)-img/p1_phy_s1_q02_1.png}{106.6}{0}
\qfigdef{p1_phy_s1_q03_1}{Score Advanced - Test 11 (09 May 2023)-img/p1_phy_s1_q03_1.png}{104.8}{0}
\qfigdef{p1_phy_s1_q05_1}{Score Advanced - Test 11 (09 May 2023)-img/p1_phy_s1_q05_1.png}{105.6}{0}
\qfigdef{p1_phy_s1_q06_1}{Score Advanced - Test 11 (09 May 2023)-img/p1_phy_s1_q06_1.png}{105.6}{0}
\qfigdef{p1_phy_s2_q07_1}{Score Advanced - Test 11 (09 May 2023)-img/p1_phy_s2_q07_1.png}{77.4}{0}
\qfigdef{p1_phy_s2_q09_1}{Score Advanced - Test 11 (09 May 2023)-img/p1_phy_s2_q09_1.png}{106.0}{0}
\qfigdef{p1_phy_s2_q09_2}{Score Advanced - Test 11 (09 May 2023)-img/p1_phy_s2_q09_2.png}{104.3}{0}
\qfigdef{p1_phy_s2_q09_3}{Score Advanced - Test 11 (09 May 2023)-img/p1_phy_s2_q09_3.png}{104.3}{0}
\qfigdef{p1_phy_s2_q09_4}{Score Advanced - Test 11 (09 May 2023)-img/p1_phy_s2_q09_4.png}{104.4}{0}
\qfigdef{p1_phy_s2_q10_1}{Score Advanced - Test 11 (09 May 2023)-img/p1_phy_s2_q10_1.png}{107.0}{0}
\qfigdef{p1_phy_s3_q02_1}{Score Advanced - Test 11 (09 May 2023)-img/p1_phy_s3_q02_1.png}{105.8}{0}
\qfigdef{p1_phy_s3_q03_1}{Score Advanced - Test 11 (09 May 2023)-img/p1_phy_s3_q03_1.png}{153.9}{0}
\qfigdef{p1_phy_s3_q04_1}{Score Advanced - Test 11 (09 May 2023)-img/p1_phy_s3_q04_1.png}{106.5}{0}
\qfigdef{p1_phy_s3_q06_1}{Score Advanced - Test 11 (09 May 2023)-img/p1_phy_s3_q06_1.png}{112.8}{0}
\qfigdef{p1_chem_s1_q03_A}{Score Advanced - Test 11 (09 May 2023)-img/p1_chem_s1_q03_A.png}{46.6}{0}
\qfigdef{p1_chem_s1_q03_B}{Score Advanced - Test 11 (09 May 2023)-img/p1_chem_s1_q03_B.png}{55.3}{0}
\qfigdef{p1_chem_s1_q03_D}{Score Advanced - Test 11 (09 May 2023)-img/p1_chem_s1_q03_D.png}{59.5}{0}
\qfigdef{p1_chem_s1_q05_1}{Score Advanced - Test 11 (09 May 2023)-img/p1_chem_s1_q05_1.png}{103.8}{0}
\qfigdef{p1_chem_s2_q09_1}{Score Advanced - Test 11 (09 May 2023)-img/p1_chem_s2_q09_1.png}{129.0}{0}
\qfigdef{p1_chem_s2_q09_B}{Score Advanced - Test 11 (09 May 2023)-img/p1_chem_s2_q09_B.png}{144.4}{0}
\qfigdef{p1_chem_s2_q09_C}{Score Advanced - Test 11 (09 May 2023)-img/p1_chem_s2_q09_C.png}{78.7}{0}
\qfigdef{p1_chem_s2_q09_D}{Score Advanced - Test 11 (09 May 2023)-img/p1_chem_s2_q09_D.png}{81.0}{0}
\qfigdef{p1_chem_s3_q03_1}{Score Advanced - Test 11 (09 May 2023)-img/p1_chem_s3_q03_1.png}{242.5}{0}
\qfigdef{p1_chem_s3_q08_1}{Score Advanced - Test 11 (09 May 2023)-img/p1_chem_s3_q08_1.png}{146.7}{0}
\raggedcolumns
\makeatletter\@beginparpenalty=10000\makeatother
\newcolumntype{L}[1]{>{\raggedright\arraybackslash}p{#1}}
\begin{document}
\begin{center}{\Large\bfseries Score Advanced --- Test 11 (09 May 2023)}\end{center}
\begin{multicols}{2}
\raggedcolumns
\papersec{1}
\partsec{1}{PHYSICS}
\sect{SECTION-I (i)}
\qn{1} Suppose two particles 1 and 2 are projected in vertical plane simultaneously.\par
\qfig{p1_phy_s1_q01_1}
Their angles of projection are 30\textdegree{} and $\theta$ , respectively, with the horizontal. Let they collide after a time t in air. Then (neglect air resistance)
\begin{qoptions}
\opt{A}{$\theta$ = sin$^{-1}$(4/5) and they will have same speed just before the collision.}
\opt{B}{$\theta$ = sin$^{-1}$(4/5) and they will have different speed just before the collision.}
\opt{C}{x $<$1280$\sqrt{3}-$960m}
\opt{D}{It is possible that the particles collide when both of them are at their highest point.}
\end{qoptions}
\qn{2} Figure shows a container filled with a liquid of density $\rho$ . Four points A, B, C and D lie on the vertices of a vertical square. Points A and C lie on a vertical line and points B and D lies on a horizontal line. Choose the correct statement(s) about the pressure at the four points.\par
\qfig{p1_phy_s1_q02_1}
\begin{qoptions}
\opt{A}{P$_\text{D}$ = P$_\text{B}$}
\opt{B}{P$_\text{A}$ $<$ P$_\text{B}$ = P$_\text{D}$ $<$ P$_\text{C}$}
\opt{C}{P$_\text{D}$ = P$_\text{B}$ = $\dfrac{\text{P}_\text{C} - \text{P}_\text{A}}{2}$}
\opt{D}{P$_\text{D}$ = P$_\text{B}$ = $\dfrac{\text{P}_\text{C} + \text{P}_\text{A}}{2}$}
\end{qoptions}
\qn{3} In the arrangement shown in figure, two positive charges, +Q each, are fixed. Mark the correct statement(s) regarding a third charged particle $-$q placed at the midpoint P that can be displaced along or perpendicular to the line connecting the charges.\par
\qfig{p1_phy_s1_q03_1}
\begin{qoptions}
\opt{A}{The particle will perform SHM for x $\ll$ a.}
\opt{B}{The particle will oscillate about P but not harmonically for any x.}
\opt{C}{The particle will perform SHM for y $\ll$ a}
\opt{D}{The particle will oscillate about P but not harmonically for y comparable to a.}
\end{qoptions}
\qn{4} A charged particle P leaves the origin with speed v = v$_0$ at some inclination with the x-axis. There is a uniform magnetic field B along the x-axis. P strikes a fixed target T on the x-axis for a minimum value of B = B$_0$. P will also strike T if.
\begin{qoptions}
\opt{A}{B = 2B$_0$ ,v = 2v$_0$}
\opt{B}{B = 2B$_0$ ,v = v$_0$}
\opt{C}{B = B$_0$ ,v = 2v$_0$}
\opt{D}{B = $\dfrac{\text{B}_0}{2}$, v = 2v$_0$}
\end{qoptions}
\qn{5} Two points monochromatic and coherent sources of light of wavelength $\lambda$ each are placed as shown in figure. The initial phase difference between the sources is zero. (D $\gg$ d). Mark the correct statement(s).\par
\qfig{p1_phy_s1_q05_1}
\begin{qoptions}
\opt{A}{If d = $\dfrac{7\lambda}{2}$ , O will be a minima.}
\opt{B}{If d = $\lambda$, only one maxima can be observed on the screen.}
\opt{C}{If d = 4.8$\lambda$, then total 5 minima would be there on the screen.}
\opt{D}{If d = $\dfrac{5\lambda}{2}$ , the intensity at O would be minimum.}
\end{qoptions}
\qn{6} One mole of an ideal monatomic gas has initial temperature T$_0$, is made to go through the cycle abca shown in the given figure. If U denotes the internal energy, then choose the correct alternative.\par
\qfig{p1_phy_s1_q06_1}
\begin{qoptions}
\opt{A}{U$_\text{c}$ $>$ U$_\text{b}$ $>$ U$_\text{a}$}
\opt{B}{U$_\text{c}$ $-$ U$_\text{b}$ = 3RT$_0$}
\opt{C}{U$_\text{c}$ $-$ U$_\text{a}$ = $\dfrac{9\text{RT}_0}{2}$}
\opt{D}{U$_\text{b}$ $-$ U$_\text{a}$ = $\dfrac{3\text{RT}_0}{2}$}
\end{qoptions}
\sect{SECTION-I (ii)}
\qn{7} In the system shown in figure, masses of the blocks are such that when system is released, the acceleration of pulley P$_1$ is a upwards and the acceleration of block 1 is a$_1$ upward. It is found that the acceleration of block 3 is same as that of 1 both in magnitude and direction. Given that a$_1$ $>$ a $>$ a$_1$/2 . Match the following.\par
\qfig{p1_phy_s2_q07_1}
{\footnotesize\setlength{\tabcolsep}{1.5pt}\noindent\begin{tabular}{|c|L{3.6cm}|c|L{2.3cm}|}
\hline
 & \multicolumn{1}{c|}{List-I} & & \multicolumn{1}{c|}{List-II}\\ \hline
(P) & Acceleration of 2 & (1) & 2a + a$_1$\\ \hline
(Q) & Acceleration of 4 & (2) & 2a $-$ a$_1$\\ \hline
(R) & Acceleration 2 w.r.t. 3 & (3) & Upward\\ \hline
(S) & Acceleration 2 w.r.t. 4 & (4) & Downward\\ \hline
\end{tabular}}\par
\begin{qoptions}
\opt{A}{P $\longrightarrow$ 2,3;Q $\longrightarrow$ 1,4;R $\longrightarrow$ 4;S $\longrightarrow$ 3}
\opt{B}{P $\longrightarrow$ 2,3;Q $\longrightarrow$ 1,2;R $\longrightarrow$ 4;S $\longrightarrow$ 3}
\opt{C}{P $\longrightarrow$ 2;Q $\longrightarrow$ 1,4;R $\longrightarrow$ 3;S $\longrightarrow$ 4}
\opt{D}{P $\longrightarrow$ 3;Q $\longrightarrow$ 1;R $\longrightarrow$ 4;S $\longrightarrow$ 3}
\end{qoptions}
\qn{8} Bodies A and B are rolling without slipping on a rough inclined plane, starting from rest. Match the following.\par
{\footnotesize\setlength{\tabcolsep}{1.5pt}\noindent\begin{tabular}{|c|L{3.1cm}|c|L{3.1cm}|}
\hline
 & \multicolumn{1}{c|}{List -I} & & \multicolumn{1}{c|}{List -II}\\ \hline
(P) & Velocity of A and B after travelling the same distance will be & (1) & Different if A and B have same size but different shape and mass\\ \hline
(Q) & Time taken by A and B to reach the bottom of the plane from the same height will be & (2) & Different if A and B have same size and same mass but having different shapes\\ \hline
(R) & Acceleration of centre of mass of A and B will be & (3) & Same if A and B have same shape and same size but of different materials\\ \hline
(S) & During rolling motion, frictional force experienced by A and B will be & (4) & Same if A and B have same size but different shape and mass\\ \hline
\end{tabular}}\par
\begin{qoptions}
\opt{A}{P $\longrightarrow$ 1,2,3;Q $\longrightarrow$ 2,3,4;R $\longrightarrow$ 1,2,3;S $\longrightarrow$ 1,2}
\opt{B}{P $\longrightarrow$ 1,2,3;Q $\longrightarrow$ 1,2,3;R $\longrightarrow$ 1,2;S $\longrightarrow$ 1,2,4}
\opt{C}{P $\longrightarrow$ 1,2,3;Q $\longrightarrow$ 2,3,4;R $\longrightarrow$ 1,2;S $\longrightarrow$ 1,2}
\opt{D}{P $\longrightarrow$ 1,2,3;Q $\longrightarrow$ 1,2,3;R $\longrightarrow$ 1,2,3;S $\longrightarrow$ 1,2}
\end{qoptions}
\qn{9} List-I represents the longitudinal vibration of metallic rods each of length 1m and each rod is rigidly clamped at two points A and B. List-II represents their frequency and wavelength of longitudinal vibration in the fundamental mode of vibration. Then match the option of List-I correctly with the options of List-II.\par
{\footnotesize\setlength{\tabcolsep}{1.5pt}\noindent\begin{tabular}{|c|L{4cm}|c|L{1.7cm}|}
\hline
 & \multicolumn{1}{c|}{List -I} & & \multicolumn{1}{c|}{List -II}\\ \hline
(P) & \qfigi{p1_phy_s2_q09_1}\par Y = 2 $\times$ 10$^{11}$N/m$^2$, $\rho$ = 8000 kg/m$^3$ & (1) & f = 40 kHz\\ \hline
(Q) & \qfigi{p1_phy_s2_q09_2}\par Y = 1.6 $\times$ 10$^{11}$N/m$^2$, $\rho$ = 2500 kg/m$^3$ & (2) & f = 12.5 kHz\\ \hline
(R) & \qfigi{p1_phy_s2_q09_3}\par Y = 2 $\times$ 10$^{11}$N/m$^2$, $\rho$ = 8000 kg/m$^3$ & (3) & f = 20 kHz\\ \hline
(S) & \qfigi{p1_phy_s2_q09_4}\par Y = 1.6 $\times$ 10$^{11}$N/m$^2$, $\rho$ = 2500 kg/m$^3$ & (4) & $\lambda$ = 0.4 m\\ \hline
 & & (5) & $\lambda$ = 0.2 m\\ \hline
\end{tabular}}\par
\begin{qoptions}
\opt{A}{P $\longrightarrow$ 1,5;Q $\longrightarrow$ 3,4;R $\longrightarrow$ 2,5;S $\longrightarrow$ 1,5}
\opt{B}{P $\longrightarrow$ 2,5;Q $\longrightarrow$ 3,5;R $\longrightarrow$ 2,4;S $\longrightarrow$ 2,5}
\opt{C}{P $\longrightarrow$ 3,4;Q $\longrightarrow$ 3,5;R $\longrightarrow$ 2,4;S $\longrightarrow$ 2,5}
\opt{D}{P $\longrightarrow$ 2,4;Q $\longrightarrow$ 3,4;R $\longrightarrow$ 2,4;S $\longrightarrow$ 1,5}
\end{qoptions}
\qn{10} In a photoelectric experimental arrangement, light of frequency f is incident on a metal target whose work function is $\phi$ = $\dfrac{\text{hf}}{3}$ as shown. In list-I, KE of photoelectron is mentioned at various locations/instants and in list-II, the corresponding values. Match the entries of list-I with the entries of list-II.\par
\qfig{p1_phy_s2_q10_1}
{\footnotesize\setlength{\tabcolsep}{1.5pt}\noindent\begin{tabular}{|c|L{4.3cm}|c|l|}
\hline
 & \multicolumn{1}{c|}{List -I} & & \multicolumn{1}{c|}{List -II}\\ \hline
(P) & Maximum KE of photoelectrons just after emission from target & (1) & Zero\\ \hline
(Q) & KE of photoelectrons just after emission from target & (2) & hf/3\\ \hline
(R) & KE of photoelectron when they are halfway between the target and collector & (3) & hf/2\\ \hline
(S) & KE of photoelectrons as they reach the collector & (4) & 2hf/3\\ \hline
\end{tabular}}\par
\begin{qoptions}
\opt{A}{P $\longrightarrow$ 4;Q $\longrightarrow$ 1,2,3,4;R $\longrightarrow$ 2,3,4;S $\longrightarrow$ 1}
\opt{B}{P $\longrightarrow$ 4;Q $\longrightarrow$ 1,4;R $\longrightarrow$ 2,3;S $\longrightarrow$ 1,3,4}
\opt{C}{P $\longrightarrow$ 1,4;Q $\longrightarrow$ 4;R $\longrightarrow$ 3,4;S $\longrightarrow$ 1,3,4}
\opt{D}{P $\longrightarrow$ 4;Q $\longrightarrow$ 1,2,3,4;R $\longrightarrow$ 1,2,3,4;S $\longrightarrow$ 1,2,3,4}
\end{qoptions}
\sect{SECTION-II}
\qn{1} A car travelling at 28 m/s has no tendency to slip (i.e. assume friction force on the car = 0) on the track of radius 200m banked at an angle $\theta$ . When the speed is increased to 35 m/s, the car is just on the point of slipping up the track. Calculate the coefficient of friction between the car and the track. (Upto two decimals) (g =10m/s$^2$.)
\qn{2} A block of mass 2kg approaches a spring of stiffness k = 80 N/m on a smooth horizontal plane with a speed of 10 m/s. Find the time in seconds (upto 2 decimals) in which kinetic energy of the block reduces to 50\% of its initial kinetic energy. Assume t = 0 when the block just touches the spring, for the first time.\par
\qfig{p1_phy_s3_q02_1}
\qn{3} A wedge of mass M=2 m$_0$ rests on a smooth horizontal plane. A small block of mass m$_0$ rests over it at left end A as shown in figure. A sharp impulse is applied on the block, due to which it starts moving to the right with velocity v$_0$ = 6 m/s . At highest point of its trajectory, the block collides with a particle of same mass m$_0$ moving vertically downwards with velocity v = 2 m/s and gets stuck with it. If the combined mass lands at the end point A of the body of mass M, calculate length $\ell$ . Neglect friction, take g = 10 m/s$^2$.\par
\qfig{p1_phy_s3_q03_1}
\qn{4} A long thin carpet is laid on the floor. One end of the carpet is bent back and then pulled backwards with constant velocity $\upsilon$ = 1 m/s , just above the part of the carpet which is still at rest on the floor. What is the minimum force (in SI system) needed to pull the moving part, if the carpet has unit length and unit mass. (Neglect all dissipative forces) (upto 1 decimal)\par
\qfig{p1_phy_s3_q04_1}
\qn{5} The volume density of the elastic deformation energy in fresh water at a depth of h = 1 km. (Take density of water (constant) =10$^3$ kg/m$^3$,g = 9.8 m/s$^2$ and bulk modulus of elasticity of water B = 2 $\times$ 10$^9$ N/m$^2$ ) is x $\times$ 10$^4$ J/m$^3$ . Find `x' (in 2 digit) (neglect atmospheric pressure).
\qn{6} Figure shows an irregular block of material of refractive index $\sqrt{2}$. A ray of light strikes the face AB as shown in the figure. After refraction it is incident on a spherical surface CD of radius of curvature 0.4m and enters a medium of refractive index 1.514 to meet PQ at E. Find the distance OE in meters (upto two places of decimal)\par
\qfig{p1_phy_s3_q06_1}
\qn{7} A photon with an energy of 4.9eV ejects photoelectrons from tungsten. When the ejected electron enters a constant magnetic field of strength B=2.5 mT at an angle of 60\textdegree{} with the field direction, the maximum pitch of the helix described by the electron is found to be 2.7 mm. Find the work function of the metal in electron volts (upto two digits). Given that specific charge of electron is 1.76$\times$10$^{11}$ C/kg
\qn{8} Consider the $\beta$ $-$ decay $^{198}$Au$\rightarrow^{198}$Hg + $\beta^-$ + $\bar{\nu}$\par
Where $^{198}$Hg represents a mercury nucleus in an excited state of energy 1.088 MeV above the ground state. What can be the maximum kinetic energy of electron emitted (upto two decimals). The atomic mass of $^{198}$Au is 197.968233u and that of $^{198}$Hg is 197.966760u. 1u = 931 MeV/c$^2$
\partsec{2}{CHEMISTRY}
\sect{SECTION-I (i)}
\qn{1} Which can form square planar complex with coordination number (CN) is 4?
\begin{qoptions}
\opt{A}{Pt(II), Pd(II) with strong and weak field ligand}
\opt{B}{Au(III) with strong and weak field ligand}
\opt{C}{Ni(II) with C.N. = 4 and strong field ligand}
\opt{D}{Hg(II) with C.N. = 4 and strong field ligand}
\end{qoptions}
\qn{2} Froth floatation:
\begin{qoptions}
\opt{A}{is physical method of separation of mineral from the gangue}
\opt{B}{is a method to concentrate the ore depending on the difference in wettability of gangue and the ore}
\opt{C}{is used for the sulphide ores}
\opt{D}{is a method in which impurities sink to the bottom}
\end{qoptions}
\qn{3} Both S$_\text{N}$1 and S$_\text{N}$2 reactions are not favourable in which of the given compounds?
\begin{qoptions}
\opt{A}{\qfigi{p1_chem_s1_q03_A}}
\opt{B}{\qfigi{p1_chem_s1_q03_B}}
\opt{C}{Ph$-$Cl}
\opt{D}{\qfigi{p1_chem_s1_q03_D}}
\end{qoptions}
\qn{4} Which of the following molecule(s) have higher resonance energy than 1,3$-$butadiene?
\begin{qoptions}
\opt{A}{Pyrrole}
\opt{B}{Thiophene}
\opt{C}{Tetrahydro Furan}
\opt{D}{1,4$-$cyclohexadiene}
\end{qoptions}
\qn{5} A mixture of fixed number of moles of ideal gases undergoes a thermodynamic process as shown by the plot. Which of the following need not be same for the process, if the composition of ideal gas mixture changes for the same initial and final temperatures? S= molar entropy of the mixture, T= temperature\par
\qfig{p1_chem_s1_q05_1}
\begin{qoptions}
\opt{A}{Change in internal energy}
\opt{B}{Change in free energy}
\opt{C}{Work done}
\opt{D}{Change in enthalpy}
\end{qoptions}
\qn{6} Given the data at 25$^\text{o}$C\par
Ag + I$^-$ $\longrightarrow$ AgI + e$^-$ ; E$^0$ = 0.152V, Ag $\longrightarrow$ Ag$^+$ + e$^-$ ; E$^0$ = $-$ 0.800V\par
What is the value of log K$_\text{sp}$ for AgI? $\left(2.303\dfrac{\text{RT}}{\text{F}} = 0.059\text{V}\right)$
\begin{qoptions}
\opt{A}{$-$ 8.12}
\opt{B}{+ 8.612}
\opt{C}{$-$ 37.83}
\opt{D}{$-$ 16.13}
\end{qoptions}
\sect{SECTION-I (ii)}
\qn{7} Match the property of the compound given in List$-$II with the compound given in List$-$I.\par
{\footnotesize\setlength{\tabcolsep}{1.5pt}\noindent\begin{tabular}{|c|l|c|L{4.3cm}|}
\hline
\multicolumn{2}{|c|}{List $-$ I} & \multicolumn{2}{c|}{List$-$ II}\\ \hline
(A) & KHCO$_3$ & (P) & contain polymeric anion\\ \hline
(B) & NaHCO$_3$ & (Q) & contain dimeric anion\\ \hline
(C) & LiHCO$_3$ & (R) & contain hydrogen bond\\ \hline
(D) & Mg(HCO$_3$)$_2$ & (S) & liberates CO$_2$ on heating\\ \hline
 & & (T) & can be prepared in solid state\\ \hline
\end{tabular}}\par
\begin{qoptions}
\opt{A}{A $\longrightarrow$ QRST;B $\longrightarrow$ PRST;C $\longrightarrow$ S;D $\longrightarrow$ ST}
\opt{B}{A $\longrightarrow$ QRST;B $\longrightarrow$ PST;C $\longrightarrow$ PS;D $\longrightarrow$ RST}
\opt{C}{A $\longrightarrow$ QST;B $\longrightarrow$ PST;C $\longrightarrow$ S;D $\longrightarrow$ RST}
\opt{D}{A $\longrightarrow$ QST;B $\longrightarrow$ PRST;C $\longrightarrow$ PS;D $\longrightarrow$ ST}
\end{qoptions}
\qn{8} Match the compound given in List$-$I with the character given in List$-$II.\par
{\footnotesize\setlength{\tabcolsep}{1.5pt}\noindent\begin{tabular}{|c|L{3.1cm}|c|L{3.1cm}|}
\hline
\multicolumn{2}{|c|}{List-I} & \multicolumn{2}{c|}{List-II}\\ \hline
(A) & Hydrogen hypo phosphate ion H$_2$P$_2$O$_6^{2-}$ & (P) & All hydrogens are ionizable\\ \hline
(B) & Hydrogen pyrophosphate ion H$_2$P$_2$O$_7^{2-}$ & (Q) & Number of acidic hydrogens present is 2\\ \hline
(C) & Pyrosulphurous acid H$_2$S$_2$O$_5$ & (R) & sp$^3$ hybridised central atom\\ \hline
(D) & Hydrogen peroxy mono sulphate ion HSO$_5^-$ & (S) & Overall non$-$planar structure\\ \hline
 & & (T) & Can act as oxidizing agent\\ \hline
\end{tabular}}\par
\begin{qoptions}
\opt{A}{A $\longrightarrow$ PQRS;B $\longrightarrow$ RST;C $\longrightarrow$ PQRS;D $\longrightarrow$ PQRST}
\opt{B}{A $\longrightarrow$ PQS;B $\longrightarrow$ RST;C $\longrightarrow$ PRS;D $\longrightarrow$ PRST}
\opt{C}{A $\longrightarrow$ PQS;B $\longrightarrow$ RS;C $\longrightarrow$ PRS;D $\longrightarrow$ PQRST}
\opt{D}{A $\longrightarrow$ PQRS;B $\longrightarrow$ PQRS;C $\longrightarrow$ PQRS;D $\longrightarrow$ PRST}
\end{qoptions}
\qn{9} Match the correct reduction method from List$-$II for the reaction in List$-$I.\par
{\footnotesize\setlength{\tabcolsep}{1.5pt}\noindent\begin{tabular}{|c|L{4cm}|c|L{2.2cm}|}
\hline
\multicolumn{2}{|c|}{List-I} & \multicolumn{2}{c|}{List-II}\\ \hline
(A) & \qfigi[0.75]{p1_chem_s2_q09_1} & (P) & Birch reduction\\ \hline
(B) & \qfigi[0.75]{p1_chem_s2_q09_B} & (Q) & Stephen's reduction\\ \hline
(C) & \qfigi[0.75]{p1_chem_s2_q09_C} & (R) & Wolff$-$Kishner reduction\\ \hline
(D) & \qfigi[0.75]{p1_chem_s2_q09_D} & (S) & Clemmensen reduction\\ \hline
\end{tabular}}\par
\begin{qoptions}
\opt{A}{A $\longrightarrow$ S;B $\longrightarrow$ R;C $\longrightarrow$ P;D $\longrightarrow$ RS}
\opt{B}{A $\longrightarrow$ R;B $\longrightarrow$ P;C $\longrightarrow$ RS;D $\longrightarrow$ S}
\opt{C}{A $\longrightarrow$ P;B $\longrightarrow$ RS;C $\longrightarrow$ S;D $\longrightarrow$ R}
\opt{D}{A $\longrightarrow$ RS;B $\longrightarrow$ S;C $\longrightarrow$ R;D $\longrightarrow$ P}
\end{qoptions}
\qn{10} Match the crystal lattice in List$-$I with the corresponding Geometrical or Physical property in List$-$II.\par
{\footnotesize\setlength{\tabcolsep}{1.5pt}\noindent\begin{tabular}{|c|L{2.5cm}|c|L{3.2cm}|}
\hline
\multicolumn{2}{|c|}{\begin{tabular}{c}List $-$ I\\(Crystal lattice)\end{tabular}} & \multicolumn{2}{c|}{\begin{tabular}{c}List $-$ II\\(Geometrical or Physical property)\end{tabular}}\\ \hline
(A) & Rock salt unit cell & (P) & Three$-$fold axis of symmetry\\ \hline
(B) & Zinc blende unit cell & (Q) & Four$-$fold axis of symmetry\\ \hline
(C) & Fluorite unit cell & (R) & Plane of symmetry\\ \hline
(D) & FCC unit cell & (S) & Anisotropy\\ \hline
 & & (T) & Number of first nearest identical species is 12\\ \hline
\end{tabular}}\par
\begin{qoptions}
\opt{A}{A $\longrightarrow$ PQST;B $\longrightarrow$ PQRS;C $\longrightarrow$ PRST;D $\longrightarrow$ PQRS}
\opt{B}{A $\longrightarrow$ PQRST;B $\longrightarrow$ PQRST;C $\longrightarrow$ PQRS;D $\longrightarrow$ PQRST}
\opt{C}{A $\longrightarrow$ PQRST;B $\longrightarrow$ QRS;C $\longrightarrow$ PQRST;D $\longrightarrow$ PQRST}
\opt{D}{A $\longrightarrow$ PQRST;B $\longrightarrow$ PQRST;C $\longrightarrow$ PQRT;D $\longrightarrow$ PQRT}
\end{qoptions}
\sect{SECTION-II}
\qn{1} A mixture contains one gram equivalent of each of the following metal hydroxides. Water is added to the mixture, boiled and filtered. How many of the hydroxides remain in the residue?\par
NaOH, KOH, Ba(OH)$_2$, TlOH, Be(OH)$_2$, Mg(OH)$_2$, Al(OH)$_3$, Sn(OH)$_2$, Zn(OH)$_2$, Sr(OH)$_2$.
\qn{2} Polyphosphates are the condensed products of tetrahedral phosphate units that contain P $-$ O $-$ P bonds. A sample of poly phosphoric acid was dissolved in water and titrated with dilute NaOH (aq). Two stoichiometric points were observed at 30.4 and 45.6 cm$^3$. The number of P $-$ O $-$ P bonds in that poly phosphoric acid is
\qn{3} \qfig{p1_chem_s3_q03_1}
Find the value of P + Q .
\qn{4} An organic compound (X) shows following characteristics\par
i) It is optically active\par
ii) Gives +ve haloform test\par
iii) Gives violet colour with neutral FeCl$_3$ solution\par
Then the minimum possible molecular weight of compound (X) is
\qn{5} How many compounds are more reactive than ethene towards electrophilic addition reaction?\par
(I) CH$_3$ $-$ CH = CH$_2$\par
(II) H $-$ C(=O) $-$ CH = CH$_2$\par
(III) Cl $-$ CH = CH$_2$\par
(IV) CH$_3$ $-$ O $-$ CH = CH$_2$\par
(V) H $-$ C $\equiv$ C $-$ H\par
(VI) Ph $-$ CH = CH$_2$\par
(VII) CH$_2$ = CH $-$ CH = CH$_2$
\qn{6} 2 mole NO and 16 gm O$_2$ were taken in a 6.25 L vessel at 27$^\text{o}$C. Calculate the change in pressure when half of the maximum amount of NO$_2$ has been produced according to reaction.\par
(Take R = $\dfrac{1}{12}$ litre.atm/mol K)\par
2NO$_{(g)}$ + O$_{2(g)}$ $\longrightarrow$ 2NO$_{2(g)}$
\qn{7} The energy of the first orbit of hydrogen atom is $-$13.6 eV. If the fourth ionization enthalpy of gaseous isolated beryllium atom is expressed as Y $\times$ 10$^{-19}$ J, what is the value of `Y'? Assume nucleus is at rest in either cases. (1 eV = 1.6 $\times$ 10$^{-19}$ J)
\qn{8} In a chemical reaction where A(g) is converting to B(g), the following observation is made. What is the time of completion for the reaction (in min)?\par
\qfig{p1_chem_s3_q08_1}
\partsec{3}{MATHEMATICS}
\sect{SECTION-I (i)}
\qn{1} $f$ : R$-$\{0\} $\longrightarrow$ R is a differentiable function such that $\displaystyle\int_1^{xy} f(t)\text{dt} = y\int_1^{x} f(t)\text{dt} + x\int_1^{y} f(t)\text{dt}$ $\forall$ x, y $\in$ R $-$ \{0\}, $f$(1) = 1.\par
Which of the following is correct?
\begin{qoptions}
\opt{A}{$\displaystyle\lim_{x\to\infty}\dfrac{f(x)}{x} = 0$}
\opt{B}{$\displaystyle\int_{1/e}^{1} f(x)\text{dx} = \dfrac{1}{e}$}
\opt{C}{Number of solutions of $f$(x) = $f^{-1}$(x) for x $>$ 0 is 1}
\opt{D}{Number of functions satisfying the given conditions can not be finite}
\end{qoptions}
\qn{2} For a given hyperbola $\dfrac{x^2}{9} - \dfrac{y^2}{4} = 1$ . Let P is any point on it such that if a tangent drawn at P is made to meet its conjugate axis in `T'. PN is the ordinate of point P and PT meets transverse axis in `T$^1$', then ( A, A$^1$ are the vertices with positive and negative x -co ordinates and C is its centre).
\begin{qoptions}
\opt{A}{CT.PN = 4}
\opt{B}{CN.CT$^1$ = 9}
\opt{C}{CN.NT$^1$ = AN.A$^1$N}
\opt{D}{CN.NT$^1$ = CN$^2$ + AC$^2$}
\end{qoptions}
\qn{3} Let a, b, c be integers and a$^2$ + b$^2$ + c$^2$ = 2 . Let A = $\begin{pmatrix} a & b & c \\ b & c & a \\ c & a & b \end{pmatrix}$ consider the system of equations A$\begin{pmatrix} x \\ y \\ z \end{pmatrix}$ = $\begin{pmatrix} 1 \\ 0 \\ 0 \end{pmatrix}$.\par
Then which of the following statement(s) is/are correct -
\begin{qoptions}
\opt{A}{The probability that the system of equations has unique solution is 1/2}
\opt{B}{The probability that the system of equations has infinitely many solutions is 1/2}
\opt{C}{If a = 0 , then the number of ordered pairs (b, c) for which the system has no solution is 2}
\opt{D}{The number of integers in the range of ab + bc + ca is 2}
\end{qoptions}
\qn{4} The plane containing the line $\dfrac{x-3}{2} = \dfrac{y-1}{4} = \dfrac{z-2}{4}$ is rotated through an angle $\dfrac{\pi}{2}$ about this line to contain the point (4, 3, 7) then which of the following statements is/are true ?
\begin{qoptions}
\opt{A}{The equation of plane before rotation is 2x + 4 y $-$ 5z = 0}
\opt{B}{The equation of plane after rotation is 2x $-$ y $-$ 5 = 0}
\opt{C}{The equation of plane before rotation is x $-$ 2y + 2z $-$ 1 = 0}
\opt{D}{The equation of plane after rotation is x + 2y + 2z = 1}
\end{qoptions}
\qn{5} For the equation $|x^2 - a| = \cos^{-1}(\cos x)$ to have some solution the value of `a' can be
\begin{qoptions}
\opt{A}{$\dfrac{1}{2}$}
\opt{B}{$\dfrac{1}{4}$}
\opt{C}{$-\dfrac{1}{2}$}
\opt{D}{$-\dfrac{1}{4}$}
\end{qoptions}
\qn{6} For $\Delta$ABC, let A$_1$B$_1$C$_1$ be the triangle formed by the excentres of $\Delta$ABC, let $\Delta$A$_2$B$_2$C$_2$ be formed by the excentres of $\Delta$A$_1$B$_1$C$_1$ and so on. Finally let $\Delta$A$_n$B$_n$C$_n$ be the n$^\text{th}$ derived triangle that it is formed by the excentres of $\Delta$A$_{n-1}$ B$_{n-1}$C$_{n-1}$ then
\begin{qoptions}
\opt{A}{$\angle$A$_1$ = 90\textdegree{} $-$ $\dfrac{A}{2}$}
\opt{B}{The circumcentre of $\Delta$ABC always coincides with the orthocentre of $\Delta$A$_1$B$_1$C$_1$}
\opt{C}{as n$\longrightarrow\infty$, $\Delta$A$_n$B$_n$C$_n$ tends to become equilateral}
\opt{D}{B$_1$C$_1$ = a cot $\dfrac{A}{2}$}
\end{qoptions}
\sect{SECTION-I (ii)}
\qn{7} Let O, O$^1$ and G be the circumcentre, orthocentre and centroid of $\Delta$ABC respectively. Now match the values in list-I and list-II.\par
{\footnotesize\setlength{\tabcolsep}{1.5pt}\noindent\begin{tabular}{|c|L{4.6cm}|c|c|}
\hline
 & \multicolumn{1}{c|}{list-I} & & \multicolumn{1}{c|}{list-II}\\ \hline
(A) & If S is any point in plane of A,B,C then $\dfrac{\left|\overrightarrow{\text{SA}} + \overrightarrow{\text{SB}} + \overrightarrow{\text{SC}}\right|}{\left|\overrightarrow{\text{SG}}\right|}$ = & (P) & 4\\ \hline
(B) & $\dfrac{\left|\overrightarrow{\text{OA}} + \overrightarrow{\text{OB}} + \overrightarrow{\text{OC}}\right|}{\left|\overrightarrow{\text{OO}'}\right|}$ = & (Q) & 3\\ \hline
(C) & $\dfrac{\left|\overrightarrow{\text{O}'\text{A}} + \overrightarrow{\text{O}'\text{B}} + \overrightarrow{\text{O}'\text{C}}\right|}{\left|\overrightarrow{\text{OO}'}\right|}$ = & (R) & 2\\ \hline
(D) & $\dfrac{\left|\overrightarrow{\text{AO}'} + \overrightarrow{\text{O}'\text{B}} + \overrightarrow{\text{O}'\text{C}}\right|}{\left|\overrightarrow{\text{AP}}\right|}$ = & (S) & 1\\ \hline
 & & (T) & 1/2\\ \hline
\end{tabular}}\par
(Where $\left|\overrightarrow{\text{AP}}\right|$ is diameter of circumcircle of $\Delta$ABC)
\begin{qoptions}
\opt{A}{A $\longrightarrow$ Q;B $\longrightarrow$ S;C $\longrightarrow$ R;D $\longrightarrow$ S}
\opt{B}{A $\longrightarrow$ Q;B $\longrightarrow$ R;C $\longrightarrow$ S;D $\longrightarrow$ S}
\opt{C}{A $\longrightarrow$ S;B $\longrightarrow$ S;C $\longrightarrow$ Q;D $\longrightarrow$ R}
\opt{D}{A $\longrightarrow$ S;B $\longrightarrow$ Q;C $\longrightarrow$ R;D $\longrightarrow$ R}
\end{qoptions}
\qn{8} \par
{\footnotesize\setlength{\tabcolsep}{1.5pt}\noindent\begin{tabular}{|c|L{5.2cm}|c|c|}
\hline
 & \multicolumn{1}{c|}{List-I} & & \multicolumn{1}{c|}{List-II}\\ \hline
(A) & If $\bar{a} = \bar{i} + \bar{j} + \bar{k}$, $\bar{b} = \bar{i} - \bar{j} + \bar{k}$, $\bar{c} = \bar{i} + 2\bar{j} - \bar{k}$ then the value of $\left|\begin{matrix} \bar{a}.\bar{a} & \bar{a}.\bar{b} & \bar{a}.\bar{c} \\ \bar{b}.\bar{a} & \bar{b}.\bar{b} & \bar{b}.\bar{c} \\ \bar{c}.\bar{a} & \bar{c}.\bar{b} & \bar{c}.\bar{c} \end{matrix}\right|$ is equal to & (P) & $\sqrt{3}$\\ \hline
(B) & If points P, Q and R have position vectors $\bar{r}_1 = 3\bar{i} - 2\bar{j} - \bar{k}$, $\bar{r}_2 = \bar{i} + 3\bar{j} + 4\bar{k}$ and $\bar{r}_3 = 2\bar{i} + \bar{j} - 2\bar{k}$ respectively relative to an origin O then the distance of P from the plane OQR is & (Q) & 28\\ \hline
(C) & A line with direction ratios (2, 1, 2) intersects the lines $\bar{r} = -\bar{j} + \lambda(\bar{i} + \bar{j} + \bar{k})$ and $\bar{r} = -\bar{i} + \mu(2\bar{i} + \bar{j} + \bar{k})$ at A and B then AB is equal to & (R) & 16\\ \hline
(D) & In a regular tetrahedron, the centres of four faces are the vertices of a smaller tetrahedron, The ratio of the volume of the smaller tetrahedron to that of the larger is $\dfrac{\text{m}}{\text{n}}$ where m and n are relatively prime positive integers. The value of m + n is & (S) & 3\\ \hline
\end{tabular}}\par
\begin{qoptions}
\opt{A}{A $\longrightarrow$ R;B $\longrightarrow$ Q;C $\longrightarrow$ S;D $\longrightarrow$ Q}
\opt{B}{A $\longrightarrow$ R;B $\longrightarrow$ S;C $\longrightarrow$ S;D $\longrightarrow$ Q}
\opt{C}{A $\longrightarrow$ S;B $\longrightarrow$ Q;C $\longrightarrow$ S;D $\longrightarrow$ Q}
\opt{D}{A $\longrightarrow$ Q;B $\longrightarrow$ R;C $\longrightarrow$ S;D $\longrightarrow$ Q}
\end{qoptions}
\qn{9} Let $f$(x) = min\{sec$^2\pi$x, 3$-\sqrt{2}$ sin$\pi$x\}, x $\in$ $\left(0, \dfrac{1}{2}\right)$\par
{\footnotesize\setlength{\tabcolsep}{1.5pt}\noindent\begin{tabular}{|c|L{4.3cm}|c|c|}
\hline
 & \multicolumn{1}{c|}{List-I} & & \multicolumn{1}{c|}{List-II}\\ \hline
(A) & Number of integers in the range of $f$(x) is & (P) & 0\\ \hline
(B) & Number of points of local extremum of $f$(x) is & (Q) & 2\\ \hline
(C) & Number of values of x where $f$(x) is an integer is & (R) & 1\\ \hline
(D) & Number of solutions of the equation $f$(x) = 1 is & (S) & 3\\ \hline
 & & (T) & 4\\ \hline
\end{tabular}}\par
\begin{qoptions}
\opt{A}{A $\longrightarrow$ S;B $\longrightarrow$ R;C $\longrightarrow$ Q;D $\longrightarrow$ P}
\opt{B}{A $\longrightarrow$ R;B $\longrightarrow$ Q;C $\longrightarrow$ Q;D $\longrightarrow$ S}
\opt{C}{A $\longrightarrow$ R;B $\longrightarrow$ R;C $\longrightarrow$ R;D $\longrightarrow$ R}
\opt{D}{A $\longrightarrow$ P;B $\longrightarrow$ S;C $\longrightarrow$ Q;D $\longrightarrow$ R}
\end{qoptions}
\qn{10} Match the following\par
{\footnotesize\setlength{\tabcolsep}{1.5pt}\noindent\begin{tabular}{|c|L{4.3cm}|c|l|}
\hline
 & \multicolumn{1}{c|}{List-I} & & \multicolumn{1}{c|}{List-II}\\ \hline
(A) & $\sqrt{\sin x}$ : $\left[0, \dfrac{\pi}{3}\right]$ $\rightarrow$ [0, 1] & (P) & one-one\\ \hline
(B) & $\dfrac{x+3}{x-1}$ : (1, $\infty$) $\rightarrow$ (1, $\infty$) & (Q) & many-one\\ \hline
(C) & sin x : $\left[-\dfrac{\pi}{2}, \dfrac{4\pi}{3}\right]$ $\rightarrow$ [$-$1, 1] & (R) & into\\ \hline
(D) & $\dfrac{x^2}{x-2}$ : (2, $\infty$) $\rightarrow$ [8, $\infty$) & (S) & onto\\ \hline
\end{tabular}}\par
\begin{qoptions}
\opt{A}{A $\longrightarrow$ P,S;B $\longrightarrow$ P,R;C $\longrightarrow$ Q,S;D $\longrightarrow$ Q,S}
\opt{B}{A $\longrightarrow$ Q,R;B $\longrightarrow$ P,S;C $\longrightarrow$ P,S;D $\longrightarrow$ Q,S}
\opt{C}{A $\longrightarrow$ P,R;B $\longrightarrow$ P,S;C $\longrightarrow$ Q,S;D $\longrightarrow$ Q,S}
\opt{D}{A $\longrightarrow$ P,R;B $\longrightarrow$ P,Q;C $\longrightarrow$ Q,S;D $\longrightarrow$ Q,S}
\end{qoptions}
\sect{SECTION-II}
\qn{1} If Z is a non-real complex number such that arg $\left(z^{1/3}\right)$ = $\dfrac{1}{2}$arg$\left(z^2 + \bar{z}z^{1/3}\right)$ then the value of $|z|$ is
\qn{2} A circle concentric with the ellipse $\dfrac{4x^2}{289} + \dfrac{y^2}{\lambda^2} = 1$ $\left(\lambda < \dfrac{17}{2}\right)$ passes through foci S$_1$, S$_2$ and cuts ellipse at point P. If area of $\Delta$PS$_1$S$_2$ is 30 sq.units then S$_1$S$_2$ =
\qn{3} Let a, b, c are roots of x$^3$ + x$^2$ $-$ 2x + 1 = 0 then the value of $\left|\begin{matrix} a+c & b & a+2b+2c \\ a+b & c & b+2a+2c \\ b+c & a & c+2a+2b \end{matrix}\right|$ is equal to
\qn{4} If a$_1$ = 1 and a$_\text{n}$= n(1 + a$_{n-1}$) $\forall$ n $\geq$ 2 then $\left(\displaystyle\lim_{n\to\infty}\left(1 + \dfrac{1}{a_1}\right)\left(1 + \dfrac{1}{a_2}\right)\ldots\ldots\left(1 + \dfrac{1}{a_n}\right)\right)\cdot\dfrac{1}{e}$ =
\qn{5} Let $f$ (x) = x$^2$ + 6x + 1 and S denotes the set of points (x, y) in the coordinate plane such that $f$(x) + $f$(y) $\leq$ 0 and $f$(x) $-$ $f$(y) $\leq$ 0. The area of S is k$\pi$ then k =
\qn{6} Let p, q$^2$ $\in$ Q, the set of rational numbers and cos$^2\dfrac{\pi}{8}$ be a root of the equation x$^2$ + px + q$^2$ = 0 then the value of [$|\cos$ p$|$ + $|\sin$ p$|$] + p + 2 $|\log_2|\text{q}||$ equals (Where[.] is G.I.F.)
\qn{7} Let x, y $\in$ R such that cos x cos y + 2sin y + 2sin x cos y = 3 then the value of tan$^2$ x + 5tan$^2$ y is
\qn{8} If $\displaystyle\lim_{n\to\infty}\left(\dfrac{{}^{3n}\text{C}_n}{{}^{2n}\text{C}_n}\right)^{\frac{1}{n}}$ = $\dfrac{a}{b}$ where a, b are relatively prime positive integers then $\dfrac{a+b+1}{a-b}$ =
\end{multicols}
\end{document}
